

\begin{frame}{La base des voyageurs}

Le schéma est résumé ci-dessous. 
\vfill

\begin{redbullet}
  \item Voyageur (\textbf{idVoyageur}, nom, prénom, ville, région)
    \item  Séjour (\textbf{idSéjour}, \textit{idVoyageur}, \textit{codeLogement}, début, fin)
    \item Logement (\textbf{code}, nom, capacité, type, lieu)
    \item Activité (\textbf{\textit{codeLogement}, codeActivité}, description)
\end{redbullet}

\vfill

En \textbf{gras}, les clés primaires, en
\textit{italiques} les clés étrangères.

\begin{remblock}{Nommage}
Totalement libre. Il est préférable d'utiliser des conventions homogènes.
\end{remblock}

\end{frame}

\begin{frame}[fragile]{Table des voyageurs et des logements}

\begin{tabular}{c|c|c|c|c}
\textbf{idVoyageur} & \textbf{nom} & \textbf{prénom} & \textbf{ville} & \textbf{région}\\ \hline 
10 & Fogg & Phileas & Ajaccio & Corse\\ \hline 
20 & Bouvier & Nicolas & Aurillac & Auvergne\\ \hline 
30 & David-Néel & Alexandra & Lhassa & Tibet\\ \hline 
40 & Stevenson & Robert Louis & Vannes & Bretagne\\ \hline 
\end{tabular}

\vfill

\begin{tabular}{c|c|c|c|c}
\textbf{code} & \textbf{nom} & \textbf{capacité} & \textbf{type} & \textbf{lieu}\\ \hline 
ca & Causses & 45 & Auberge & Cévennes\\ \hline 
ge & Génépi & 134 & Hôtel & Alpes\\ \hline 
pi & U Pinzutu & 10 & Gîte & Corse\\ \hline 
ta & Tabriz & 34 & Hôtel & Bretagne\\ \hline 
\end{tabular}
\end{frame}


\begin{frame}[fragile]{Table des séjours}

\begin{center}
\begin{tabular}{c|c|c|c|c}
\textbf{idSéjour} & \textbf{\textit{idVoyageur}} & \textbf{\textit{codeLogement}} & \textbf{début} & \textbf{fin}\\ \hline 
1 & 10 & pi & 20 & 20\\ \hline 
2 & 20 & ta & 21 & 22\\ \hline 
3 & 30 & ge & 2 & 3\\ \hline 
4 & 20 & pi & 19 & 23\\ \hline 
5 & 20 & ge & 22 & 24\\ \hline 
6 & 10 & pi & 10 & 12\\ \hline 
7 & 30 & ca & 13 & 18\\ \hline 
\end{tabular}
\end{center}

\vfill

\begin{remblock}{Rôle des clés étrangères}
Chaque nuplet référence \alert{un} voyageur et \alert{un} logement.
Mais un voyageur ou un logement peut être
référencé plusieurs fois.
\end{remblock}
\end{frame}



\begin{frame}[fragile]{Table des activités}
\begin{center}
\begin{tabular}{c|c|c}
\textbf{\textit{codeLog.}} & \textbf{codeAct.} & \textbf{description}\\ \hline 
ca & Randonnée & Sorties d’une journée en groupe\\ \hline 
ge & Piscine & Nage loisir non encadrée\\ \hline 
ge & Ski & Sur piste uniquement\\ \hline 
pi & Plongée & Baptèmes et préparation des brevets\\ \hline 
pi & Voile & Pratique du dériveur et du catamaran\\ \hline 
\end{tabular}
\end{center}


\vfill

\begin{remblock}{Clé primaire et étrangère }
Une clé étrangère peut être une partie d'une clé primaire.
\end{remblock}
\end{frame}

%%%%%%ù

\begin{frame}{À retenir}

Comprendre un schéma relationnel.

\vfill
\begin{redbullet}
  \item Identifier la clé primaire de chaque relation (il doit \alert{toujours} y en avoir une)
 \item Identifier les clés étrangères; comprendre quelle clé primaire elles référencent.
 \item Comprendre comment les nuplets sont liés transitivement les uns aux autres.
\end{redbullet}

\vfill

Ces deux bases sont accessibles en ligne: \url{http://deptfod.cnam.fr/bd/tp}
\end{frame}

